\documentclass[a4paper,11pt]{article}
\usepackage[T1]{fontenc}
\usepackage[utf8]{inputenc}
\usepackage[left=2cm,right=2cm,top=2.5cm,bottom=2.5cm]{geometry}
\setlength{\headheight}{14pt}
\usepackage{xcolor}
\usepackage{mdframed}
\usepackage{parskip}
\usepackage{fancyhdr}
\usepackage{hyperref}

\definecolor{usercolor}{RGB}{30, 100, 200}
\definecolor{agentcolor}{RGB}{30, 150, 80}
\definecolor{doccolor}{RGB}{200, 90, 10}
\definecolor{userbg}{RGB}{235, 244, 255}
\definecolor{agentbg}{RGB}{235, 250, 238}
\definecolor{docbg}{RGB}{255, 243, 220}
\definecolor{answerbg}{RGB}{250, 250, 235}

\newmdenv[backgroundcolor=userbg,linecolor=usercolor,linewidth=1.5pt,
  leftmargin=0pt,rightmargin=80pt,innerleftmargin=10pt,innerrightmargin=10pt,
  innertopmargin=8pt,innerbottommargin=8pt,roundcorner=5pt,
  skipabove=6pt,skipbelow=3pt]{userbubble}

\newmdenv[backgroundcolor=agentbg,linecolor=agentcolor,linewidth=1.5pt,
  leftmargin=80pt,rightmargin=0pt,innerleftmargin=10pt,innerrightmargin=10pt,
  innertopmargin=8pt,innerbottommargin=8pt,roundcorner=5pt,
  skipabove=3pt,skipbelow=3pt]{agentbubble}

\newmdenv[backgroundcolor=docbg,linecolor=doccolor,linewidth=1pt,
  leftmargin=80pt,rightmargin=0pt,innerleftmargin=10pt,innerrightmargin=10pt,
  innertopmargin=6pt,innerbottommargin=6pt,roundcorner=4pt,
  skipabove=2pt,skipbelow=3pt]{docenv}

\newmdenv[backgroundcolor=answerbg,linecolor=gray,linewidth=0.5pt,
  leftmargin=80pt,rightmargin=0pt,innerleftmargin=10pt,innerrightmargin=10pt,
  innertopmargin=5pt,innerbottommargin=5pt,roundcorner=4pt,
  skipabove=2pt,skipbelow=6pt]{refanswerenv}

\pagestyle{fancy}
\fancyhf{}
\renewcommand{\headrulewidth}{0.3pt}
\cfoot{\thepage}

\lhead{\textcolor{gray}{\small RAG Conversation Dataset}}
\rhead{\textcolor{gray}{\small Conv.\,\#2\ --\ govt0611}}
\begin{document}

\begin{center}
{\LARGE\bfseries RAG Conversation #8}
\par\medskip
{\small\textcolor{gray}{Domain:\ \texttt{mt-rag-ibmcloud-elser-512-100-20240502}\quad|\quad Status:\ \textbf{accepted}}}
\end{center}
\vspace{0.4em}
\noindent\textcolor{gray}{\rule{\linewidth}{0.6pt}}
\vspace{0.4em}
\vspace{0.5em}
\noindent{\small\textcolor{usercolor}{\textbf{$\triangleright$\enspace User --- Turn~1}}}\hfill{\tiny\textcolor{gray}{Factoid \$\\textbackslash{}cdot\$ ANSWERABLE}}
\begin{userbubble}
does IBM offer document databases?
\end{userbubble}
\noindent\hfill{\small\textcolor{agentcolor}{\textbf{Agent\enspace$\triangleleft$}}}
\begin{agentbubble}
Yes, IBM offers [IBM Cloudant](https://cloud.ibm.com/docs/Cloudant?topic=Cloudant-documents) which stores JSON objects.
\end{agentbubble}
\noindent\hfill{\footnotesize\textcolor{doccolor}{\textbf{$\hookrightarrow$\enspace Retrieved documents~(3)}}}
\begin{docenv}
{\footnotesize\textcolor{doccolor}{\textbf{Doc~1}}}\hfill{\tiny\texttt{score:~18.445}}\ {\tiny\textcolor{gray}{\{'relevant': \{'8163b85d-d338-4ed1-8ebf-cd065ac3dc12': \{'value': 'yes', 'timestamp': 1724868041\}\}\}}}\\[3pt]
{\small\textit{(sans titre)}}\\[2pt]
{\tiny\url{https://cloud.ibm.com/docs/Cloudant?topic=Cloudant-combine-ibm-services}}\\[4pt]
{\footnotesize 
Summary 

Today, we combined two IBM Cloud services to optimize cost and user experience: IBM Cloudant as a document store and query engine and Databases for Redis as a content cache. Cached documents can be retrieved more quickly and more cheaply, but the tradeoff is that your application might be showing old data to your users for a time.

If you followed this tutorial, you must deprovision your resources to stop incurring charges. You can deprovision from the terraform directory on your terminal by typing the following command, terraform destroy --auto-approve.}
\end{docenv}
\begin{docenv}
{\footnotesize\textcolor{doccolor}{\textbf{Doc~2}}}\hfill{\tiny\texttt{score:~18.217}}\ {\tiny\textcolor{gray}{\{'relevant': \{'8163b85d-d338-4ed1-8ebf-cd065ac3dc12': \{'value': 'yes', 'timestamp': 1724857361\}\}\}}}\\[3pt]
{\small\textit{(sans titre)}}\\[2pt]
{\tiny\url{https://cloud.ibm.com/docs/Cloudant?topic=Cloudant-databases}}\\[4pt]
{\footnotesize 
Database overview 

IBM® Cloudant® for IBM Cloud® databases contain JSON objects. These JSON objects are called [documents](https://cloud.ibm.com/docs/Cloudant?topic=Cloudant-documentsdocuments).

All documents must be contained in a database. For more information, see [partitioned databases](https://cloud.ibm.com/docs/Cloudant?topic=Cloudant-databasespartitioned-databases-database).

The [Grouping related documents together in IBM Cloudant](https://cloud.ibm.com/docs/Cloudant?topic=Cloudant-grouping-related-documents-together-in-ibm-cloudantgrouping-related-documents-together-in-ibm-cloudant) guide provides an example of how documents for an e-commerce application might be used within an IBM Cloudant database.



 Partitioned databases 

IBM Cloudant supports two types of databases:



* Partitioned
* Nonpartitioned



A partitioned database offers significant query performance and cost advantages but requires you to specify a logical partitioning of your data. The partitioning is specified as part of each document's ID. A partitioned database provides both global and partition queries. Partition queries target queries at a single, given document partition, meaning they need to process less data to return results. Therefore, partition queries offer significant performance advantages, and also often provide cost advantages over global queries. Global queries target the entire database, which leads to extra complexity, slower performance, and increased cost, but offers results that draw from all data.

Alternatively, a nonpartitioned database might be created. This type of database can be less complex to work with since no partitioning scheme needs to be defined, but you can create only global secondary indexes.

IBM Cloudant strongly encourages you to use a partitioned database for best long-term database performance where the data model allows for logical partitioning of documents.

The partitioning type of a database is set at database creation time. When you create a database, use the partitioned query string parameter to set whether the database is partitioned. The default for partitioned is false, maintaining compatibility with an earlier version.}
\end{docenv}
\begin{docenv}
{\footnotesize\textcolor{doccolor}{\textbf{Doc~3}}}\hfill{\tiny\texttt{score:~18.014}}\ {\tiny\textcolor{gray}{\{'relevant': \{'8163b85d-d338-4ed1-8ebf-cd065ac3dc12': \{'value': 'no', 'timestamp': 1724857709\}\}\}}}\\[3pt]
{\small\textit{(sans titre)}}\\[2pt]
{\tiny\url{https://cloud.ibm.com/docs/databases-for-cassandra}}\\[4pt]
{\footnotesize 
IBM Cloud® Databases for DataStax docs 

IBM Cloud® Databases for DataStax is a scale-out NoSQL database built on Apache Cassandra, designed for high-availability and workload flexibility.

 Developer tools 

[API \& SDK reference](https://cloud.ibm.com/apidocs/cloud-databases-api/cloud-databases-api-v5introduction)[CLI reference](https://cloud.ibm.com/docs/databases-cli-plugin)[Terraform reference](https://registry.terraform.io/providers/IBM-Cloud/ibm/latest/docs/resources/database)

 Recommended content 

[Deprecation of IBM Cloud® Databases for DataStax IBM Cloud® Databases for DataStax is deprecated and no longer supported as of 30 June 2024.](https://cloud.ibm.com/docs/databases-for-cassandra?topic=databases-for-cassandra-deprecation)[Setting the Admin Password The Databases for DataStax service is provisioned with an admin user. You must set the admin password before you use it to connect.](https://cloud.ibm.com/docs/databases-for-cassandra?topic=databases-for-cassandra-admin-password\&interface=cli)[Creating and managing users, roles, and privileges IBM Cloud® Databases for DataStax uses a system of roles to manage database permissions. Roles are used to give a single user or a group of users a set of privileges.](https://cloud.ibm.com/docs/databases-for-cassandra?topic=databases-for-cassandra-user-management\&interface=cli)

 Learn more 

[IBM Developer<br><br>Visit IBM Developer for technical articles, code patterns, tutorials, and more.<br><br>!}
\end{docenv}
\begin{refanswerenv}
{\footnotesize\textit{\textcolor{gray}{\textbf{Reference answer:}\enspace Yes, IBM offers [IBM Cloudant](https://cloud.ibm.com/docs/Cloudant?topic=Cloudant-documentsdocuments).}}}
\end{refanswerenv}
\vspace{0.5em}
\noindent{\small\textcolor{usercolor}{\textbf{$\triangleright$\enspace User --- Turn~2}}}\hfill{\tiny\textcolor{gray}{Follow-up \$\\textbackslash{}cdot\$ Composite, Factoid \$\\textbackslash{}cdot\$ ANSWERABLE}}
\begin{userbubble}
So it can store any random JSON object or I need to specify fields in advances?
\end{userbubble}
\noindent\hfill{\small\textcolor{agentcolor}{\textbf{Agent\enspace$\triangleleft$}}}
\begin{agentbubble}
Yes, IBM Cloudant can store any JSON object subject to some caveats listed below
1. All documents must have two fields:
    * A unique \_id field. The \_id field is detailed in the next section.
    * A \_rev field. The \_rev field is a revision identifier, and is essential to the IBM Cloudant replication protocol.
2. Field names that begin with the underscore character (\_) are reserved in IBM Cloudant. This rule means that you can't normally have your own field names that begin with an underscore. For example, the field example would be accepted, but the field \_example would result in a doc\_validation error message.

\end{agentbubble}
\noindent\hfill{\footnotesize\textcolor{doccolor}{\textbf{$\hookrightarrow$\enspace Retrieved documents~(2)}}}
\begin{docenv}
{\footnotesize\textcolor{doccolor}{\textbf{Doc~1}}}\hfill{\tiny\texttt{score:~31.144}}\ {\tiny\textcolor{gray}{\{'relevant': \{'8163b85d-d338-4ed1-8ebf-cd065ac3dc12': \{'value': 'yes', 'timestamp': 1724858280\}\}\}}}\\[3pt]
{\small\textit{(sans titre)}}\\[2pt]
{\tiny\url{https://cloud.ibm.com/docs/Cloudant?topic=Cloudant-documents}}\\[4pt]
{\footnotesize 
Overview 

Documents are [JSON objects](https://en.wikipedia.org/wiki/JSONData\_types.2C\_syntax\_and\_example). Documents are also containers for your data, and are the basis of the IBM® Cloudant® for IBM Cloud® database.

If you're using an [IBM Cloudant service on IBM Cloud®](https://cloud.ibm.com/docs/Cloudant?topic=Cloudant-ibm-cloud-publicibm-cloud-public), documents are limited to a maximum size of 1 MB. Exceeding this limit causes a [413 error](https://cloud.ibm.com/apidocs/cloudantlist-of-http-codes).

IBM Cloudant uses an [eventually consistent](https://cloud.ibm.com/docs/Cloudant?topic=Cloudant-cap-theoremcap-theorem) model for data. If you use the eventually consistent model, it's possible, under some conditions, to retrieve older document content. For example, older content is retrieved when your application writes or updates a document that is followed immediately by a read of the same document.

In other words, your application would see the document content as it was before the write or update occurred. For more information about this model, see the topic on [Consistency](https://cloud.ibm.com/docs/Cloudant?topic=Cloudant-cap-theoremcap-theorem).



 Document fields 

All documents must have two fields:



* A unique \_id field. The \_id field is detailed in the next section.
* A \_rev field. The \_rev field is a revision identifier, and is [essential to the IBM Cloudant replication protocol](https://cloud.ibm.com/docs/Cloudant?topic=Cloudant-document-versioning-and-mvccdocument-versioning-and-mvcc).



In addition to these two mandatory fields, documents can generally contain any other content that can be described by using JSON, subject to some caveats detailed in the following sections.



 Document IDs}
\end{docenv}
\begin{docenv}
{\footnotesize\textcolor{doccolor}{\textbf{Doc~2}}}\hfill{\tiny\texttt{score:~20.475}}\ {\tiny\textcolor{gray}{\{'relevant': \{'8163b85d-d338-4ed1-8ebf-cd065ac3dc12': \{'value': 'yes', 'timestamp': 1724858283\}\}\}}}\\[3pt]
{\small\textit{(sans titre)}}\\[2pt]
{\tiny\url{https://cloud.ibm.com/docs/Cloudant?topic=Cloudant-documents}}\\[4pt]
{\footnotesize 
Field name restrictions 

Field names that begin with the underscore character (\_) are reserved in IBM Cloudant. This rule means that you can't normally have your own field names that begin with an underscore. For example, the field example would be accepted, but the field \_example would result in a doc\_validation error message.

See an example of a JSON document that attempts to create a field with an underscore prefix:

\{
"\_top\_level\_field\_name": "some data"
\}

See an error message that is returned when you attempt to create a field with an underscore prefix:

\{
"error": "doc\_validation",
"reason": "Bad special document member: \_top\_level\_field\_name"
\}

However, if the field name is for an object that is nested within the document, you can use an underscore prefix for the field name.

See an example of a JSON document that attempts to create a field with an underscore prefix, nested within an object:

\{
"another\_top\_level\_field\_name": "some data",
"another\_field": \{
"\_lower\_level\_field\_name": "some more data"
\}
\}

See an example success message (abbreviated) returned when a nested field with an underscore prefix is created:

\{
"ok": true,
"id": "2",
"rev": "1-9ce...8d4"
\}







 Quorum - writing and reading data 

In a distributed system, it's possible that a request might take some time to complete. A "quorum" mechanism is used to help determine when a request, such as a write or read, completes successfully.

For more information about quorum settings and their implications on dedicated IBM Cloudant systems, contact IBM Cloudant support.





 Time to live}
\end{docenv}
\begin{refanswerenv}
{\footnotesize\textit{\textcolor{gray}{\textbf{Reference answer:}\enspace Yes, IBM Cloudant can store any JSON object.}}}
\end{refanswerenv}
\vspace{0.5em}
\noindent{\small\textcolor{usercolor}{\textbf{$\triangleright$\enspace User --- Turn~3}}}\hfill{\tiny\textcolor{gray}{Factoid \$\\textbackslash{}cdot\$ Follow-up \$\\textbackslash{}cdot\$ ANSWERABLE}}
\begin{userbubble}
What if I want to store an image or pdf with a document? 
\end{userbubble}
\noindent\hfill{\small\textcolor{agentcolor}{\textbf{Agent\enspace$\triangleleft$}}}
\begin{agentbubble}
IBM Cloudant has support for storing attachments alongside documents, a long-standing feature it inherits from CouchDB. So you can easily add an image or PDF as an attachment to your document but you must consider a few things before you use attachments in IBM Cloudant today, especially if you’re looking at larger assets.

1. IBM Cloudant is expensive as a block store.
2. IBM Cloudant’s internal implementation is not efficient in handling large amounts of binary data. So, slow and expensive. IBM Cloudant is acceptable for small assets and occasional use. As a rule, if you need to store binary data alongside IBM Cloudant documents, it’s better to use a separate solution more suited for this purpose. You need store only the attachment metadata in the IBM Cloudant document.
\end{agentbubble}
\noindent\hfill{\footnotesize\textcolor{doccolor}{\textbf{$\hookrightarrow$\enspace Retrieved documents~(1)}}}
\begin{docenv}
{\footnotesize\textcolor{doccolor}{\textbf{Doc~1}}}\hfill{\tiny\texttt{score:~17.237}}\ {\tiny\textcolor{gray}{\{'relevant': \{'8163b85d-d338-4ed1-8ebf-cd065ac3dc12': \{'value': 'yes', 'timestamp': 1724859734\}\}\}}}\\[3pt]
{\small\textit{(sans titre)}}\\[2pt]
{\tiny\url{https://cloud.ibm.com/docs/Cloudant?topic=Cloudant-data-modeling}}\\[4pt]
{\footnotesize 
IBM Cloudant has support for storing attachments alongside documents, a long-standing feature it inherits from CouchDB. If you use IBM Cloudant as a backend for a web application, you can also store small icons and other static assets such as CSS and JavaScript files with the data.

You must consider a few things before you use attachments in IBM Cloudant today, especially if you’re looking at larger assets such as images and videos:



1. IBM Cloudant is expensive as a block store.
2. IBM Cloudant’s internal implementation is not efficient in handling large amounts of binary data.



So, slow and expensive.

IBM Cloudant is acceptable for small assets and occasional use. As a rule, if you need to store binary data alongside IBM Cloudant documents, it’s better to use a separate solution more suited for this purpose. You need store only the attachment metadata in the IBM Cloudant document. Yes, that means you need to write some extra code to upload the attachment to a suitable block store of your choice. Verify that it succeeded before you store the token or URL to the attachment in the IBM Cloudant document.

Your databases are smaller, cheaper, faster, and easier to replicate. For more information, see the following websites:



* IBM Cloudant docs on [attachments](https://cloud.ibm.com/docs/Cloudant?topic=Cloudant-how-to-use-attachments)
* Detaching IBM Cloudant attachments to [Object Storage](https://medium.com/codait/detaching-cloudant-attachments-to-object-storage-with-serverless-functions-99b8c3c77925)







 Fewer databases are better than many 

If you can, limit the number of databases per IBM Cloudant account to 500 or fewer. While this particular number is not magic (IBM Cloudant can safely handle more), several use cases exist that are adversely affected by large numbers of databases in an account.

The replicator scheduler has a limited number of simultaneous replication jobs that it is prepared to run. As the number of databases grows, the replication latency is likely to increase if you try to replicate everything contained in an account.}
\end{docenv}
\begin{refanswerenv}
{\footnotesize\textit{\textcolor{gray}{\textbf{Reference answer:}\enspace You need to store only the attachment metadata in the IBM Cloudant document. Yes, that means you need to write some extra code to upload the attachment to a suitable block store of your choice. Verify that it succeeded before you store the token or URL to the attachment in the IBM Cloudant document.}}}
\end{refanswerenv}
\vspace{0.5em}
\noindent{\small\textcolor{usercolor}{\textbf{$\triangleright$\enspace User --- Turn~4}}}\hfill{\tiny\textcolor{gray}{Factoid \$\\textbackslash{}cdot\$ Follow-up \$\\textbackslash{}cdot\$ PARTIAL}}
\begin{userbubble}
Is there a limit on the file size? 
\end{userbubble}
\noindent\hfill{\small\textcolor{agentcolor}{\textbf{Agent\enspace$\triangleleft$}}}
\begin{agentbubble}
Yes, documents are limited to a maximum size of 1 MB. Exceeding this limit causes a [413 error](https://cloud.ibm.com/apidocs/cloudantlist-of-http-codes).
\end{agentbubble}
\noindent\hfill{\footnotesize\textcolor{doccolor}{\textbf{$\hookrightarrow$\enspace Retrieved documents~(1)}}}
\begin{docenv}
{\footnotesize\textcolor{doccolor}{\textbf{Doc~1}}}\hfill{\tiny\texttt{score:~41.154}}\ {\tiny\textcolor{gray}{\{'relevant': \{'8163b85d-d338-4ed1-8ebf-cd065ac3dc12': \{'value': 'yes', 'timestamp': 1724867561\}\}\}}}\\[3pt]
{\small\textit{(sans titre)}}\\[2pt]
{\tiny\url{https://cloud.ibm.com/docs/Cloudant?topic=Cloudant-documents}}\\[4pt]
{\footnotesize 
Overview 

Documents are [JSON objects](https://en.wikipedia.org/wiki/JSONData\_types.2C\_syntax\_and\_example). Documents are also containers for your data, and are the basis of the IBM® Cloudant® for IBM Cloud® database.

If you're using an [IBM Cloudant service on IBM Cloud®](https://cloud.ibm.com/docs/Cloudant?topic=Cloudant-ibm-cloud-publicibm-cloud-public), documents are limited to a maximum size of 1 MB. Exceeding this limit causes a [413 error](https://cloud.ibm.com/apidocs/cloudantlist-of-http-codes).

IBM Cloudant uses an [eventually consistent](https://cloud.ibm.com/docs/Cloudant?topic=Cloudant-cap-theoremcap-theorem) model for data. If you use the eventually consistent model, it's possible, under some conditions, to retrieve older document content. For example, older content is retrieved when your application writes or updates a document that is followed immediately by a read of the same document.

In other words, your application would see the document content as it was before the write or update occurred. For more information about this model, see the topic on [Consistency](https://cloud.ibm.com/docs/Cloudant?topic=Cloudant-cap-theoremcap-theorem).



 Document fields 

All documents must have two fields:



* A unique \_id field. The \_id field is detailed in the next section.
* A \_rev field. The \_rev field is a revision identifier, and is [essential to the IBM Cloudant replication protocol](https://cloud.ibm.com/docs/Cloudant?topic=Cloudant-document-versioning-and-mvccdocument-versioning-and-mvcc).



In addition to these two mandatory fields, documents can generally contain any other content that can be described by using JSON, subject to some caveats detailed in the following sections.



 Document IDs}
\end{docenv}
\vspace{0.5em}
\noindent{\small\textcolor{usercolor}{\textbf{$\triangleright$\enspace User --- Turn~5}}}\hfill{\tiny\textcolor{gray}{Explanation \$\\textbackslash{}cdot\$ Clarification \$\\textbackslash{}cdot\$ PARTIAL}}
\begin{userbubble}
Not documents, I was asking for attachments?
\end{userbubble}
\noindent\hfill{\small\textcolor{agentcolor}{\textbf{Agent\enspace$\triangleleft$}}}
\begin{agentbubble}
I am sorry I couldn't find any explicit limits on attachment size but as stated earlier IBM Cloudant isn't best suited for handling large amounts of binary data. Thus it’s better to use a separate solution more suited for this purpose and only store the attachment metadata in the IBM Cloudant document.
\end{agentbubble}
\noindent\hfill{\footnotesize\textcolor{doccolor}{\textbf{$\hookrightarrow$\enspace Retrieved documents~(1)}}}
\begin{docenv}
{\footnotesize\textcolor{doccolor}{\textbf{Doc~1}}}\hfill{\tiny\texttt{score:~45.349}}\ {\tiny\textcolor{gray}{\{'relevant': \{'8163b85d-d338-4ed1-8ebf-cd065ac3dc12': \{'value': 'yes', 'timestamp': 1724866360\}\}\}}}\\[3pt]
{\small\textit{(sans titre)}}\\[2pt]
{\tiny\url{https://cloud.ibm.com/docs/Cloudant?topic=Cloudant-data-modeling}}\\[4pt]
{\footnotesize 
IBM Cloudant has support for storing attachments alongside documents, a long-standing feature it inherits from CouchDB. If you use IBM Cloudant as a backend for a web application, you can also store small icons and other static assets such as CSS and JavaScript files with the data.

You must consider a few things before you use attachments in IBM Cloudant today, especially if you’re looking at larger assets such as images and videos:



1. IBM Cloudant is expensive as a block store.
2. IBM Cloudant’s internal implementation is not efficient in handling large amounts of binary data.



So, slow and expensive.

IBM Cloudant is acceptable for small assets and occasional use. As a rule, if you need to store binary data alongside IBM Cloudant documents, it’s better to use a separate solution more suited for this purpose. You need store only the attachment metadata in the IBM Cloudant document. Yes, that means you need to write some extra code to upload the attachment to a suitable block store of your choice. Verify that it succeeded before you store the token or URL to the attachment in the IBM Cloudant document.

Your databases are smaller, cheaper, faster, and easier to replicate. For more information, see the following websites:



* IBM Cloudant docs on [attachments](https://cloud.ibm.com/docs/Cloudant?topic=Cloudant-how-to-use-attachments)
* Detaching IBM Cloudant attachments to [Object Storage](https://medium.com/codait/detaching-cloudant-attachments-to-object-storage-with-serverless-functions-99b8c3c77925)







 Fewer databases are better than many 

If you can, limit the number of databases per IBM Cloudant account to 500 or fewer. While this particular number is not magic (IBM Cloudant can safely handle more), several use cases exist that are adversely affected by large numbers of databases in an account.

The replicator scheduler has a limited number of simultaneous replication jobs that it is prepared to run. As the number of databases grows, the replication latency is likely to increase if you try to replicate everything contained in an account.}
\end{docenv}
\begin{refanswerenv}
{\footnotesize\textit{\textcolor{gray}{\textbf{Reference answer:}\enspace assistant: IBM Cloudant has support for storing attachments alongside documents, a long-standing feature it inherits from CouchDB. So you can easily add an image or PDF as an attachment to your document but you must consider a few things before you use attachments in IBM Cloudant today, especially if you’re looking at larger assets. 1. IBM Cloudant is expensive as a block store. 2. IBM Cloudant’s internal implementation is not efficient in handling large amounts of binary data. So, slow and expensive. IBM Cloudant is acceptable for small assets and occasional use. As a rule, if you need to store binary data alongside IBM Cloudant documents, it’s better to use a separate solution more suited for this purpose. You need store only the attachment metadata in the IBM Cloudant document.}}}
\end{refanswerenv}
\vspace{0.5em}
\noindent{\small\textcolor{usercolor}{\textbf{$\triangleright$\enspace User --- Turn~6}}}\hfill{\tiny\textcolor{gray}{How-To \$\\textbackslash{}cdot\$ Follow-up \$\\textbackslash{}cdot\$ ANSWERABLE}}
\begin{userbubble}
Ok. How do I get started? I just look through cloud catelog for IBM Cloudant? 
\end{userbubble}
\noindent\hfill{\small\textcolor{agentcolor}{\textbf{Agent\enspace$\triangleleft$}}}
\begin{agentbubble}
It's quite easy to get setup with IBM Cloudant instance. Let me walk you through it.

1. Log in to your IBM Cloud account, and click Create resource. 

2. Type Cloudant in the Search bar and click to open it. 

3. Select an offering and an environment. 

4. Type an instance name. 

5. Select your plan. 

6. To create the service, click Create:
After you click Create, the system displays a message to say that the instance is being provisioned, which returns you to the Resource list. From the Resource list, you see that the status for your instance is, Provision in progress. 

7. After you create an instance, the status changes to Active.  
\end{agentbubble}
\noindent\hfill{\footnotesize\textcolor{doccolor}{\textbf{$\hookrightarrow$\enspace Retrieved documents~(1)}}}
\begin{docenv}
{\footnotesize\textcolor{doccolor}{\textbf{Doc~1}}}\hfill{\tiny\texttt{score:~18.284}}\ {\tiny\textcolor{gray}{\{'relevant': \{'8163b85d-d338-4ed1-8ebf-cd065ac3dc12': \{'value': 'yes', 'timestamp': 1724867012\}\}\}}}\\[3pt]
{\small\textit{(sans titre)}}\\[2pt]
{\tiny\url{https://cloud.ibm.com/docs/Cloudant?topic=Cloudant-getting-started-with-cloudant}}\\[4pt]
{\footnotesize 
Getting started with IBM Cloudant 

The IBM® Cloudant® for IBM Cloud® Getting started tutorial demonstrates how to use the IBM Cloud® dashboard to create an IBM Cloudant service instance and obtain service credentials to connect to it. Finally, it guides you through the creation of a simple, locally hosted web application that uses your IBM Cloudant database.



 Objectives 



* Create a service instance.
* Create an IBM Cloudant service credential.







 Step 1: Creating a service instance 



1. Log in to your IBM Cloud account, and click Create resource.

Zoom

![IBM Cloud Dashboard, which includes Build tile, Monitor your resources tile, Create and deploy an application tile, API Connect tile, Integrate Watson with anything tile, and Watson starter kits tile.](https://cloud.ibm.com/docs-content/v1/content/522c6f62358b063f921283400a601c3f9bc66f08/Cloudant//tutorials/images/img0001.png)

Figure 1. IBM Cloud Dashboard

The IBM Cloud Dashboard can be found at: [https://cloud.ibm.com/](https://cloud.ibm.com/). After you authenticate with your username and password, you're presented with the IBM Cloud Dashboard.
2. Type Cloudant in the Search bar and click to open it.
3. Select an offering and an environment.
4. Type an instance name.

Zoom

![Create the IBM Cloudant service name and credentials.](https://cloud.ibm.com/docs-content/v1/content/522c6f62358b063f921283400a601c3f9bc66f08/Cloudant/tutorials/images/img0005b.png)

Figure 2. IBM Cloudant service name and credentials

(In this example, the instance name is Cloudant-o7.) Verify that the resource group and authentication methods are correct. Add a tag if you like.}
\end{docenv}
\begin{refanswerenv}
{\footnotesize\textit{\textcolor{gray}{\textbf{Reference answer:}\enspace assistant: IBM Cloudant is available in the [IBM Cloud](https://cloud.ibm.com/docs/Cloudant?topic=Cloudant-documents) and the [IBM Cloudant](https://cloud.ibm.com/docs/Cloudant?topic=Cloudant-documents) cloud catelog.}}}
\end{refanswerenv}
\vspace{0.5em}
\noindent{\small\textcolor{usercolor}{\textbf{$\triangleright$\enspace User --- Turn~7}}}\hfill{\tiny\textcolor{gray}{Comparative \$\\textbackslash{}cdot\$ Follow-up \$\\textbackslash{}cdot\$ ANSWERABLE}}
\begin{userbubble}
What are difference between lite and standard plan?
\end{userbubble}
\noindent\hfill{\small\textcolor{agentcolor}{\textbf{Agent\enspace$\triangleleft$}}}
\begin{agentbubble}
The Lite plan is free but has a fixed amount of throughput capacity. The paid Standard plan offers configurable provisioned throughput capacity and data storage pricing that scales as your application requirements change.  

Here are the couple of specific differences
- In the case of a Lite plan,  provisioned throughput capacity is fixed at 20 reads per second, 10 writes per second, and 5 global queries per second, and data storage is capped at 1 GB. A Standard plan starts with a provisioned throughput capacity of 100 reads per second, 50 writes per second, and 5 global queries per second with 20 GB of data storage with an easy way to increase them in small increments.  
- You're limited to one IBM Cloudant Lite plan instance per IBM Cloud account.
\end{agentbubble}
\noindent\hfill{\footnotesize\textcolor{doccolor}{\textbf{$\hookrightarrow$\enspace Retrieved documents~(3)}}}
\begin{docenv}
{\footnotesize\textcolor{doccolor}{\textbf{Doc~1}}}\hfill{\tiny\texttt{score:~19.979}}\ {\tiny\textcolor{gray}{\{'relevant': \{'8163b85d-d338-4ed1-8ebf-cd065ac3dc12': \{'value': 'yes', 'timestamp': 1724943892\}\}\}}}\\[3pt]
{\small\textit{(sans titre)}}\\[2pt]
{\tiny\url{https://cloud.ibm.com/docs/services/Cloudant?topic=Cloudant-ibm-cloud-public}}\\[4pt]
{\footnotesize 
Plans and provisioning 

IBM® Cloudant® for IBM Cloud® Standard is IBM Cloudant's most feature-rich offering, receiving updates and new features first. Pricing is based on provisioned throughput capacity that is allocated and data storage that is used, making it suitable for any required load.

The free [Lite plan](https://cloud.ibm.com/docs/services/Cloudant?topic=Cloudant-ibm-cloud-publiclite-plan) includes a fixed amount of throughput capacity and data for development and evaluation purposes. The paid [Standard plan](https://cloud.ibm.com/docs/services/Cloudant?topic=Cloudant-ibm-cloud-publicstandard-plan) offers configurable provisioned throughput capacity and data storage pricing that scales as your application requirements change. An optional [Dedicated Hardware plan](https://cloud.ibm.com/docs/services/Cloudant?topic=Cloudant-ibm-cloud-publicdedicated-hardware-plan) is also available for an extra monthly fee to run one or more of your Standard plan instances on a dedicated hardware environment. The dedicated hardware environment is for your sole use. If a Dedicated Hardware plan instance is provisioned within a US location, you can optionally select a [HIPAA](https://en.wikipedia.org/wiki/Health\_Insurance\_Portability\_and\_Accountability\_Act) -compliant configuration.



 Plans 

You can select which plan to use when you [provision your IBM Cloudant service instance](https://cloud.ibm.com/docs/services/Cloudant?topic=Cloudant-ibm-cloud-publicprovisioning-a-cloudant-nosql-db-instance-on-ibm-cloud). The available plans include Lite and Standard. When you select a plan, Capacity displays, and the Cost estimator shows the monthly charge for the selected plan. By default, the [Lite plan](https://cloud.ibm.com/docs/services/Cloudant?topic=Cloudant-ibm-cloud-publiclite-plan) is selected.



 Lite plan 

The Lite plan is free, and is designed for development and evaluation purposes.}
\end{docenv}
\begin{docenv}
{\footnotesize\textcolor{doccolor}{\textbf{Doc~2}}}\hfill{\tiny\texttt{score:~18.774}}\ {\tiny\textcolor{gray}{\{'relevant': \{'8163b85d-d338-4ed1-8ebf-cd065ac3dc12': \{'value': 'yes', 'timestamp': 1724943887\}\}\}}}\\[3pt]
{\small\textit{(sans titre)}}\\[2pt]
{\tiny\url{https://cloud.ibm.com/docs/Cloudant?topic=Cloudant-ibm-cloud-public}}\\[4pt]
{\footnotesize 
The available plans include Lite and Standard. When you select a plan, Capacity displays, and the Cost estimator shows the monthly charge for the selected plan. By default, the [Lite plan](https://cloud.ibm.com/docs/Cloudant?topic=Cloudant-ibm-cloud-publiclite-plan) is selected.



 Lite plan 

The Lite plan is free, and is designed for development and evaluation purposes. All of IBM Cloudant's functions are included, but Lite plan instances have a fixed amount of provisioned throughput capacity and data storage. The provisioned throughput capacity is fixed at 20 reads per second, 10 writes per second, and 5 global queries per second, and data storage is capped at 1 GB.

Storage usage is checked daily. If you exceed your 1-GB storage limit, requests to the IBM Cloudant instance receive a 402 status code with the following error message, Account exceeded its data usage quota. An upgrade to a paid plan is required. A banner also appears on the IBM Cloudant Dashboard. You can still read and delete data. However, to write new data, you have two options. First, you can upgrade to a paid [Standard plan](https://cloud.ibm.com/docs/Cloudant?topic=Cloudant-ibm-cloud-publicstandard-plan), which removes the write limitation immediately. Alternatively, you can delete data so that your total storage falls under the 1-GB limit and wait until the next daily storage check runs for your instance to allow writes again.

If you want to store more than one GB of data, or be able to scale provisioned throughput capacity, move to the [Standard plan](https://cloud.ibm.com/docs/Cloudant?topic=Cloudant-ibm-cloud-publicstandard-plan).

You're limited to one IBM Cloudant Lite plan instance per IBM Cloud account. If you already have one Lite plan instance, you can't create a second Lite plan instance, or change a Standard plan instance to a Lite plan. If you try, you see the following message. You can have only one instance of a Lite plan per service.}
\end{docenv}
\begin{docenv}
{\footnotesize\textcolor{doccolor}{\textbf{Doc~3}}}\hfill{\tiny\texttt{score:~20.238}}\ {\tiny\textcolor{gray}{\{'relevant': \{'8163b85d-d338-4ed1-8ebf-cd065ac3dc12': \{'value': 'yes', 'timestamp': 1724943889\}\}\}}}\\[3pt]
{\small\textit{(sans titre)}}\\[2pt]
{\tiny\url{https://cloud.ibm.com/docs/Cloudant?topic=Cloudant-ibm-cloud-public}}\\[4pt]
{\footnotesize 
If you want to store more than one GB of data, or be able to scale provisioned throughput capacity, move to the [Standard plan](https://cloud.ibm.com/docs/Cloudant?topic=Cloudant-ibm-cloud-publicstandard-plan).

You're limited to one IBM Cloudant Lite plan instance per IBM Cloud account. If you already have one Lite plan instance, you can't create a second Lite plan instance, or change a Standard plan instance to a Lite plan. If you try, you see the following message. You can have only one instance of a Lite plan per service. To create a new instance, either delete your existing Lite plan instance or select a paid plan.





 Standard plan 

The IBM Cloudant Standard plan is available to all paid IBM Cloud® accounts, either as pay-as-you-go or subscription, and scales to meet the needs of your application. The Standard plan is priced based on two factors: the provisioned throughput capacity that is allocated, and the amount of data that is stored in the instance.

Pricing is pro-rated hourly with a starting provisioned throughput capacity of 100 reads per second, 50 writes per second, and 5 global queries per second. This rate is equal to a starting cost of USD \$0.105 per hour. You can toggle the provisioned throughput capacity up or down by using the user interface or API. Toggle in increments of 100 reads per second, 50 writes per second, and 5 global queries per second. Costs are calculated for the provisioned throughput capacity that is allocated and not on the metered volume of requests. The Standard plan includes 20 GB of data storage. If you store more than 20 GB, you're charged a defined cost per GB per hour.

Refer to the IBM Cloud Pricing Calculator in the dashboard for pricing at different capacities and currencies, and the [pricing](https://cloud.ibm.com/docs/Cloudant?topic=Cloudant-pricingpricing) information for examples to estimate costs.





 Dedicated Hardware plan}
\end{docenv}
\begin{refanswerenv}
{\footnotesize\textit{\textcolor{gray}{\textbf{Reference answer:}\enspace assistant: The Lite plan is free, and is designed for development and evaluation purposes. The paid [Standard plan](https://cloud.ibm.com/docs/services/Cloudant?topic=Cloudant-ibm-cloud-publicstandard-plan) offers configurable provisioned throughput capacity and data storage pricing that scales as your application requirements change.}}}
\end{refanswerenv}
\vspace{0.5em}
\noindent{\small\textcolor{usercolor}{\textbf{$\triangleright$\enspace User --- Turn~8}}}\hfill{\tiny\textcolor{gray}{Opinion \$\\textbackslash{}cdot\$ Follow-up \$\\textbackslash{}cdot\$ ANSWERABLE}}
\begin{userbubble}
So does that mean if I am just exploring things for now, I should stick to to Lite plan? 
\end{userbubble}
\noindent\hfill{\small\textcolor{agentcolor}{\textbf{Agent\enspace$\triangleleft$}}}
\begin{agentbubble}
As long as you don't require more that 1 GB of data storage and comfortable with a fixed throughput capacity of 20 reads per second, 10 writes per second, and 5 global queries per second, you can start with a Lite plan. Additionally, if you change you mind later, you can always upgrade from a Lite to Standard plan.
\end{agentbubble}
\noindent\hfill{\footnotesize\textcolor{doccolor}{\textbf{$\hookrightarrow$\enspace Retrieved documents~(3)}}}
\begin{docenv}
{\footnotesize\textcolor{doccolor}{\textbf{Doc~1}}}\hfill{\tiny\texttt{score:~20.626}}\ {\tiny\textcolor{gray}{\{'relevant': \{'8163b85d-d338-4ed1-8ebf-cd065ac3dc12': \{'value': 'yes', 'timestamp': 1724944075\}\}\}}}\\[3pt]
{\small\textit{(sans titre)}}\\[2pt]
{\tiny\url{https://cloud.ibm.com/docs/services/Cloudant?topic=Cloudant-ibm-cloud-public}}\\[4pt]
{\footnotesize 
Plans and provisioning 

IBM® Cloudant® for IBM Cloud® Standard is IBM Cloudant's most feature-rich offering, receiving updates and new features first. Pricing is based on provisioned throughput capacity that is allocated and data storage that is used, making it suitable for any required load.

The free [Lite plan](https://cloud.ibm.com/docs/services/Cloudant?topic=Cloudant-ibm-cloud-publiclite-plan) includes a fixed amount of throughput capacity and data for development and evaluation purposes. The paid [Standard plan](https://cloud.ibm.com/docs/services/Cloudant?topic=Cloudant-ibm-cloud-publicstandard-plan) offers configurable provisioned throughput capacity and data storage pricing that scales as your application requirements change. An optional [Dedicated Hardware plan](https://cloud.ibm.com/docs/services/Cloudant?topic=Cloudant-ibm-cloud-publicdedicated-hardware-plan) is also available for an extra monthly fee to run one or more of your Standard plan instances on a dedicated hardware environment. The dedicated hardware environment is for your sole use. If a Dedicated Hardware plan instance is provisioned within a US location, you can optionally select a [HIPAA](https://en.wikipedia.org/wiki/Health\_Insurance\_Portability\_and\_Accountability\_Act) -compliant configuration.



 Plans 

You can select which plan to use when you [provision your IBM Cloudant service instance](https://cloud.ibm.com/docs/services/Cloudant?topic=Cloudant-ibm-cloud-publicprovisioning-a-cloudant-nosql-db-instance-on-ibm-cloud). The available plans include Lite and Standard. When you select a plan, Capacity displays, and the Cost estimator shows the monthly charge for the selected plan. By default, the [Lite plan](https://cloud.ibm.com/docs/services/Cloudant?topic=Cloudant-ibm-cloud-publiclite-plan) is selected.



 Lite plan 

The Lite plan is free, and is designed for development and evaluation purposes.}
\end{docenv}
\begin{docenv}
{\footnotesize\textcolor{doccolor}{\textbf{Doc~2}}}\hfill{\tiny\texttt{score:~19.954}}\ {\tiny\textcolor{gray}{\{'relevant': \{'8163b85d-d338-4ed1-8ebf-cd065ac3dc12': \{'value': 'yes', 'timestamp': 1724944146\}\}\}}}\\[3pt]
{\small\textit{(sans titre)}}\\[2pt]
{\tiny\url{https://cloud.ibm.com/docs/Cloudant?topic=Cloudant-ibm-cloud-public}}\\[4pt]
{\footnotesize 
The available plans include Lite and Standard. When you select a plan, Capacity displays, and the Cost estimator shows the monthly charge for the selected plan. By default, the [Lite plan](https://cloud.ibm.com/docs/Cloudant?topic=Cloudant-ibm-cloud-publiclite-plan) is selected.



 Lite plan 

The Lite plan is free, and is designed for development and evaluation purposes. All of IBM Cloudant's functions are included, but Lite plan instances have a fixed amount of provisioned throughput capacity and data storage. The provisioned throughput capacity is fixed at 20 reads per second, 10 writes per second, and 5 global queries per second, and data storage is capped at 1 GB.

Storage usage is checked daily. If you exceed your 1-GB storage limit, requests to the IBM Cloudant instance receive a 402 status code with the following error message, Account exceeded its data usage quota. An upgrade to a paid plan is required. A banner also appears on the IBM Cloudant Dashboard. You can still read and delete data. However, to write new data, you have two options. First, you can upgrade to a paid [Standard plan](https://cloud.ibm.com/docs/Cloudant?topic=Cloudant-ibm-cloud-publicstandard-plan), which removes the write limitation immediately. Alternatively, you can delete data so that your total storage falls under the 1-GB limit and wait until the next daily storage check runs for your instance to allow writes again.

If you want to store more than one GB of data, or be able to scale provisioned throughput capacity, move to the [Standard plan](https://cloud.ibm.com/docs/Cloudant?topic=Cloudant-ibm-cloud-publicstandard-plan).

You're limited to one IBM Cloudant Lite plan instance per IBM Cloud account. If you already have one Lite plan instance, you can't create a second Lite plan instance, or change a Standard plan instance to a Lite plan. If you try, you see the following message. You can have only one instance of a Lite plan per service.}
\end{docenv}
\begin{docenv}
{\footnotesize\textcolor{doccolor}{\textbf{Doc~3}}}\hfill{\tiny\texttt{score:~21.095}}\ {\tiny\textcolor{gray}{\{'relevant': \{'8163b85d-d338-4ed1-8ebf-cd065ac3dc12': \{'value': 'yes', 'timestamp': 1724944148\}\}\}}}\\[3pt]
{\small\textit{(sans titre)}}\\[2pt]
{\tiny\url{https://cloud.ibm.com/docs/services/Cloudant?topic=Cloudant-ibm-cloud-public}}\\[4pt]
{\footnotesize 
If you want to store more than one GB of data, or be able to scale provisioned throughput capacity, move to the [Standard plan](https://cloud.ibm.com/docs/services/Cloudant?topic=Cloudant-ibm-cloud-publicstandard-plan).

You're limited to one IBM Cloudant Lite plan instance per IBM Cloud account. If you already have one Lite plan instance, you can't create a second Lite plan instance, or change a Standard plan instance to a Lite plan. If you try, you see the following message. You can have only one instance of a Lite plan per service. To create a new instance, either delete your existing Lite plan instance or select a paid plan.





 Standard plan 

The IBM Cloudant Standard plan is available to all paid IBM Cloud® accounts, either as pay-as-you-go or subscription, and scales to meet the needs of your application. The Standard plan is priced based on two factors: the provisioned throughput capacity that is allocated, and the amount of data that is stored in the instance.

Pricing is pro-rated hourly with a starting provisioned throughput capacity of 100 reads per second, 50 writes per second, and 5 global queries per second. This rate is equal to a starting cost of USD \$0.105 per hour. You can toggle the provisioned throughput capacity up or down by using the user interface or API. Toggle in increments of 100 reads per second, 50 writes per second, and 5 global queries per second. Costs are calculated for the provisioned throughput capacity that is allocated and not on the metered volume of requests. The Standard plan includes 20 GB of data storage. If you store more than 20 GB, you're charged a defined cost per GB per hour.

Refer to the IBM Cloud Pricing Calculator in the dashboard for pricing at different capacities and currencies, and the [pricing](https://cloud.ibm.com/docs/Cloudant?topic=Cloudant-pricingpricing) information for examples to estimate costs.





 Dedicated Hardware plan}
\end{docenv}
\begin{refanswerenv}
{\footnotesize\textit{\textcolor{gray}{\textbf{Reference answer:}\enspace assistant: The Lite plan is free but has a fixed amount of throughput capacity. The paid Standard plan offers configurable provisioned throughput capacity and data storage pricing that scales as your application requirements change. Here are the couple of specific differences - In the case of a Lite plan, provisioned throughput capacity is fixed at 20 reads per second, 10 writes per second, and 5 global queries per second, and data storage is capped at 1 GB. A Standard plan starts with a provisioned throughput capacity of 100 reads per second, 50 writes per second, and 5 global queries per second with 20 GB of data storage with an easy way to increase them in small increments. - You're limited to one IBM Cloudant Lite plan instance per IBM Cloud account.}}}
\end{refanswerenv}
\vspace{0.5em}
\noindent{\small\textcolor{usercolor}{\textbf{$\triangleright$\enspace User --- Turn~9}}}\hfill{\tiny\textcolor{gray}{Factoid \$\\textbackslash{}cdot\$ Follow-up \$\\textbackslash{}cdot\$ ANSWERABLE}}
\begin{userbubble}
Great! One last thing, where can I find python SDK for pushing and fetching data? 
\end{userbubble}
\noindent\hfill{\small\textcolor{agentcolor}{\textbf{Agent\enspace$\triangleleft$}}}
\begin{agentbubble}
The IBM Cloudant SDK for Python is available at https://pypi.org/project/ibmcloudant/.
\end{agentbubble}
\noindent\hfill{\footnotesize\textcolor{doccolor}{\textbf{$\hookrightarrow$\enspace Retrieved documents~(2)}}}
\begin{docenv}
{\footnotesize\textcolor{doccolor}{\textbf{Doc~1}}}\hfill{\tiny\texttt{score:~17.944}}\ {\tiny\textcolor{gray}{\{'relevant': \{'8163b85d-d338-4ed1-8ebf-cd065ac3dc12': \{'value': 'yes', 'timestamp': 1724944924\}\}\}}}\\[3pt]
{\small\textit{(sans titre)}}\\[2pt]
{\tiny\url{https://cloud.ibm.com/docs/Cloudant?topic=Cloudant-client-libraries}}\\[4pt]
{\footnotesize 
The IBM Cloudant SDK for Python library is the official IBM Cloudant library for Python.

Install the [IBM Cloudant SDK for Python](https://pypi.org/project/ibmcloudant/) library by running pip or easy\_install as shown in the following examples:

pip install --upgrade "ibmcloudant>=0.0.27"

or

easy\_install --upgrade "ibmcloudant>=0.0.27"

For more information, see the [python.org](https://www.python.org/about/) website.



 Library for Python 



* [IBM Cloudant SDK for Python](https://github.com/IBM/cloudant-python-sdk)









 Go 

The IBM Cloudant SDK for Go library is the official IBM Cloudant library for Go.

Install the [IBM Cloudant SDK for Go](https://pkg.go.dev/mod/github.com/IBM/cloudant-go-sdk) library by running the following command:

go get -u github.com/IBM/cloudant-go-sdk/cloudantv1



 Library for Go 



* [IBM Cloudant SDK for Go](https://github.com/ibm/cloudant-go-sdk)









 Useful tools 

You can use the following tools with IBM Cloudant.



 Supported tools 

Supported tools are maintained and supported by IBM Cloudant.



 couchbackup 

A tool that you use from the command line to back up an IBM Cloudant or CouchDB database to a text file.

To install couchbackup, run the following command by using npm:

npm install -g @cloudant/couchbackup

For more information, see [couchbackup](https://github.com/cloudant/couchbackup).







 Unsupported tools 

Unsupported tools are not maintained or supported by IBM Cloudant.



 cURL 

A tool that you use from the command line to transfer data.

For more information, see [curl](https://curl.haxx.se/).}
\end{docenv}
\begin{docenv}
{\footnotesize\textcolor{doccolor}{\textbf{Doc~2}}}\hfill{\tiny\texttt{score:~17.767}}\ {\tiny\textcolor{gray}{\{'relevant': \{'8163b85d-d338-4ed1-8ebf-cd065ac3dc12': \{'value': 'yes', 'timestamp': 1724944921\}\}\}}}\\[3pt]
{\small\textit{(sans titre)}}\\[2pt]
{\tiny\url{https://cloud.ibm.com/docs/Cloudant?topic=Cloudant-client-libraries}}\\[4pt]
{\footnotesize 
Client libraries 

Client libraries are the tools that you use to develop your own applications to work with IBM® Cloudant® for IBM Cloud® databases.

The following client libraries are formally supported by IBM Cloudant: Java™, Node.js, Python, and Go.

A supported library is one where you can contact IBM Cloudant if you come across a specific, reproducible problem in the current version of the library.



 Java 

The IBM Cloudant SDK for Java™ is the official IBM Cloudant library for Java™.

To install the IBM Cloudant SDK for the Java™ library, see [Installation](https://github.com/ibm/cloudant-java-sdkinstallation) about installing the library by adding it as a dependency to your Maven or Gradle builds. You can also see details and examples of how to use the library in the guide.



 Library for Java 



* [IBM Cloudant SDK for Java™](https://github.com/ibm/cloudant-java-sdk)









 Node.js 

The IBM Cloudant SDK for the Node.js library is the official IBM Cloudant library for Node.js.

Install the [IBM Cloudant SDK for Node.js](https://www.npmjs.com/package/@ibm-cloud/cloudant) library by running the following command:

npm install @ibm-cloud/cloudant



 Library for Node.js 



* [IBM Cloudant SDK for Node.js](https://github.com/ibm/cloudant-node-sdk)









 Python 

The IBM Cloudant SDK for Python library is the official IBM Cloudant library for Python.

Install the [IBM Cloudant SDK for Python](https://pypi.org/project/ibmcloudant/) library by running pip or easy\_install as shown in the following examples:

pip install --upgrade "ibmcloudant>=0.0.27"

or

easy\_install --upgrade "ibmcloudant>=0.0.27"}
\end{docenv}
\begin{refanswerenv}
{\footnotesize\textit{\textcolor{gray}{\textbf{Reference answer:}\enspace assistant: The IBM Cloudant SDK for Python is available at https://pypi.org/project/ibmcloudant/.}}}
\end{refanswerenv}
\end{document}
